% \iffalse meta-comment
%
% Copyright (C) 2017- by Yanshuo Chu <yanshuoc@gmail.com>
%
% This file may be distributed and/or modified under the
% conditions of the LaTeX Project Public License, either version 1.3a
% of this license or (at your option) any later version.
% The latest version of this license is in:
%
% http://www.latex-project.org/lppl.txt
%
% and version 1.3a or later is part of all distributions of LaTeX
% version 2004/10/01 or later.
%
% \fi
%
% ^^A ======================================================================
% ^^A hithesis 用户手册
% ^^A
% ^^A 只放面向用户的说明（模板介绍 / 安装 / 使用说明），不含任何实现代码，
% ^^A 也不参与 docstrip：hithesis.ins 从不读这个文件。
% ^^A 由 hithesis.dtx 的 driver 通过 \DocInput 引入，渲染成手册的前半部分。
% ^^A ======================================================================
%
% \GetFileInfo{\jobname.dtx}
%
%
% \def\indexname{索引}
% \def\glossaryname{修改记录}
% \IndexPrologue{\section{\indexname}}
% \GlossaryPrologue{\section{\glossaryname}}
%
% \title{\bfseries\color{violet}\hithesis：哈尔滨工业大学学位论文模板}
% \author{{\fangsong 初砚硕\thanks{\href{mailto:yanshuoc@gmail.com}{yanshuoc@gmail.com}}，孙忠豪\thanks{\href{mailto:syvshc@foxmail.com}{syvshc@foxmail.com}}，石佳乐\thanks{\href{mailto:mingzzx@outlook.com}{mingzzx@outlook.com}}，李云鹏\thanks{\href{mailto:ln3dqagt03gnqxkqt0@163.com}{ln3dqagt03gnqxkqt0@163.com}}，吴俊豪\thanks{\href{mailto:junhaowu_hit@163.com}{junhaowu\_hit@163.com}}，支卓\thanks{\href{mailto:xeonacid@hit.edu.cn}{xeonacid@hit.edu.cn}}}}
% \date{v\fileversion\ (\filedate)}
% \maketitle\thispagestyle{empty}
% \begin{abstract}\noindent
% 该宏包为哈尔滨工业大学一校三区本、硕、博中英文毕业论文、开题和中期报告以及博后出站报
% 告模板。
% \end{abstract}
%
% \vskip2cm
% \def\abstractname{免责声明}
% \begin{abstract}
% \noindent
% \begin{enumerate}
% \item 本模板的发布遵守 \LaTeX\ Project Public License，使用前请认真阅读协议内
%   容。
% \item 本模板为作者根据\hit 教务处颁发的\UGR ，\hit 研究生院颁发的\PGR 编写而成
% ，为方便\hit 学生撰写毕业论文使用。
% \item \hit 教务处和研究生院只提供毕业论文写作指南，不提供官方模板（包括MS word
% ），也不会授权第三方模板为官方模板\footnote{目前深圳校区的本科模板基于\hithesis 修改得来。}，所以此模板仅为写作指南的参考实现，不保证格
% 式审查老师不提意见。任何由于使用本模板而引起的论文格式审查问题均与本模板作者无
% 关。
% \item 任何个人或组织以本模板为基础进行修改、扩展而生成的新的专用模板，请严格遵
%   守 \LaTeX\ Project Public License 协议。由于违犯协议而引起的任何纠纷争端均与
%   本模板作者无关。
% \end{enumerate}
% \end{abstract}
%
%
% \clearpage
% \pagestyle{fancy}
% \begin{multicols}{2}[
%   \setlength{\columnseprule}{.4pt}
%   \setlength{\columnsep}{18pt}]
%   \tableofcontents
% \end{multicols}
% \clearpage
%
% \section{模板介绍}
% \hithesis\ (\textbf{H}arbin \textbf{I}nstitute of \textbf{T}echnology \LaTeX\
% \textbf{Thesis} Template) 是为了帮助\hit 毕业生撰写毕业论文而编写
% 的 \LaTeX\ 论文模板。
%
% 本文档将尽量完整的介绍模板的使用方法，如有不清楚之处可以参考示例文档或者根据
% 第~\ref{sec:howtoask} 节说明提问，有兴趣者都可以参与完善此手册，也非常欢迎对代
% 码的贡献。
%
% \subsection*{文档记号的说明}
% 本文采用不同字体表示不同内容。无衬线字体表示宏包名称，如
% \pkg{xeCJK} 宏包、\cls{hithesisbook} 文档类等；等宽字体表示代码或
% 文件名，如 \cs{hitsetup} 命令、\env{abstract} 环境、\TeX{} 文档
% \file{thesis.tex} 等；带有尖括号的楷体（或西文斜体）表示命令参数，
% 如 \option{模板选项}、\option{English title} 等。在使用时，参数两侧
% 的尖括号不必输入。示例代码进行了语法高亮处理，以方便阅读。
% 在用户手册中，带有\textcolor{violet}{紫色}侧边线的为 \LaTeX{} 代码，
% 而带有\textcolor{purple}{红色}侧边线的则为命令行代码，请注意区分。
%
% \note[注意：]{模板的作用在于减少论文写作过程中格式调整的时间。前提是遵守模板的
% 用法，否则即便用了 \hithesis\ 也难以保证输出的论文符合学校规范。}
%
%
% \section{安装}
% \label{sec:installation}
% ^^A 未来
% ^^A \hithesis\ 将已经包含在主要的 \TeX\ 发行版中，一般不需要安装，可以利用发行版自
% ^^A 带更新工具自动更新。阅读文档可以使用以下命令：
% ^^A \begin{shell}
% ^^A $ texdoc hithesis
% ^^A \end{shell} 
% \textbf{目前还没有将最新版的 \hithesis\ 上传到 CTAN，所以不要使用 \texttt{texdoc} 来打开文档。}
%
% 可以在相关 QQ 群内下载编译好的模板文件和文档，但是文档还没有更新完毕，有一些内容是过时的，有一些更新的功能没有添加到文档中，等写好文档了就传 CTAN，又挖了个坑。
%
% 如果要使用开发版，需自己下载，\hithesis\ 相关链接是 
% \begin{itemize}
% \item github：\href{https://github.com/hithesis/hithesis}
% {https://github.com/hithesis/hithesis}
% \item github dev 分支：\href{https://github.com/hithesis/hithesis/tree/dev}
% {https://github.com/hithesis/hithesis/tree/dev}
% \end{itemize}
% 其中 dev 分支中的为最新开发版, 有可能有还没解决的 bug （目前来看好像不太可能）, 如果没有额外描述, 请尽量使用 dev 分支的模板
% \subsection{模板的组成}
% 下表列出了 \hithesis\ 的主要文件及其功能介绍：
%
% \begin{longtable}{lp{8cm}}
% \toprule
% {\heiti 文件（夹）} & {\heiti 功能描述}\\\midrule
% \endfirsthead
% \midrule
% {\heiti 文件（夹）} & {\heiti 功能描述}\\\midrule
% \endhead
% \endfoot
% \endlastfoot
% \file{hithesis.ins} & \textsc{DocStrip} 驱动文件（开发用） \\
% \file{hithesis.dtx} & \textsc{DocStrip} 主源文件（开发用）\\
% \file{hithesis-eps.dtx} & EPS 图像数据（校徽、封面、示例插图等，从主源剥离）\\
% \file{dtx-style.sty} & \textsc{DocStrip} 样式文件（开发用）\\\midrule
% \file{hithesisbook.cls} & 模板类文件\\
% \file{hithesisbook.cfg} & 模板配置文件\\
% \file{hithesisart.cls} & 开题和中期报告模板类文件 \\
% \file{hithesisart.cfg} & 开题和中期报告模板配置文件 \\
% \file{hithesisartplus.cls} & 深圳校区博士中期报告模板类文件 \\
% \file{hithesis.bst} & 参考文献样式文件\\
% \file{hitszthesis.bst} & 深圳校区参考文献样式文件 \\\midrule
% \file{hithesis.ist} & 索引样式文件\\\midrule
% \file{examples/} & 模板示例文件夹 \\
% \texttt{├─ }\file{hitart/} & 开题和中期报告文件夹\\
% \texttt{│\phantom{─ }├─ }\file{reports/} & 开题和中期报告文件夹 \\
% \texttt{│\phantom{─ }└─ }\file{reportplus/} & 深圳校区博士中期报告文件夹 \\
% \texttt{└─ }\file{hitbook/} & 毕业论文文件夹\\
% \texttt{\phantom{├─ }├─ }\file{chinese/} & 中文毕业论文文件夹 \\
% \texttt{\phantom{├─ }└─ }\file{english/} & 英文毕业论文文件夹 \\\midrule
% \file{Makefile} & Makefile\\
% \file{latexmkrc} & latexmk 配置文件 \\
% \file{README.md} & Readme\\
% \file{\textbf{hithesis.pdf}} & 用户手册（本文档）\\\bottomrule
% \end{longtable}
%
% 几点说明：
% \begin{itemize}
% \item 所有的 \file{hithesis*.cls}、\file{hithesis*.cfg} 以及 \file{hithesis.bst}、
%   \file{hitszthesis.bst}、\file{hithesis.ist} 都可由 \file{hithesis.ins} 和 \file{hithesis.dtx} 生成。
% \item 使用前阅读文档： \file{hithesis.pdf} 。
% \item 使用 \file{hithesis.ins} 生成样式文件后， \file{examples/} 文件夹中的 \file{chinese/}，\file{english/} ，\file{reports/} 与 \file{reportplus/} 文件夹均可以单独使用。
% \end{itemize}
%
% \subsection{代码组织}
% \label{sec:modules}
% \changes{v3.2a}{0000/00/00}{类文件内部按职责分区重写，用户接口不变}
% v3.2a 把类文件的实现按职责分成了若干区块（选项、字体、元信息、封面、章节标题、参考文献、
% 页眉页脚、目录等），每个区块只管一件事。这些区块在生成 \file{*.cls} 时按固定顺序拼接，
% 最终仍是三个独立的类文件，用户接口与此前完全一致。
%
% 分区只影响源码的组织方式，用户不会接触到。维护者若要了解各区块的职责划分与新增区块的
% 步骤，见仓库根目录的 \file{CONTRIBUTING.md}。
%
% \subsection{生成模板}
% \label{sec:generate-cls}
% \note[提示：]{若使用 \TeX 发行版自带的 \hithesis\ 或 Github/OSChina
% 上发布的版本，可忽略此节，直接阅读第 \ref{sec:generate-thesis} 节。若下载
% CTAN 包或者 Github/OSChina 开发代码，请阅读本节了解生成模板文件的步骤。}
%
% 模板解压缩后生成文件夹 \file{hithesis-vX.Y.Z} 或 \file{hithesis-vX.YZ}\footnote{\texttt{vX.Y.Z} 与 \texttt{vX.YZ} 为版本号。}，
% 其中包括：模板源文件（\file{hithesis.ins} 和 \file{hithesis.dtx}），参考文献
% 样式 \file{hithesis.bst}，示例文档
% （\file{thesis.tex}\footnote{自 v3.2a 起，原 \file{hithesis.sty} 的代码已按职责并入
% 类文件内部，不再单独生成该文件，用户接口完全不变。}，\file{front/}、\file{body/}、
% \file{back/} 和 \file{figures/}）。在使用之前需要先生成模
% 板文件和配置文件（具体命令细节请参考 \file{README.md} 和 \file{Makefile}）：
%
% \begin{shell}
% $ cd hithesis-vX.Y.Z
% # 生成 hithesisbook.cls 和 hithesisbook.cfg
% $ latex hithesis.ins
%
% # 下面的命令用来生成用户手册，可以不执行
% $ xelatex hithesis.dtx
% $ makeindex -s gind.ist -o hithesis.ind hithesis.idx
% $ makeindex -s gglo.ist -o hithesis.gls hithesis.glo
% $ xelatex hithesis.dtx
% $ xelatex hithesis.dtx  % 生成说明文档 hithesis.pdf
% \end{shell}
%
% \subsection{生成论文}
% \label{sec:generate-thesis}
% 本节介绍几种常见的生成论文的方法。用户可根据自己的情况选择。
%
% \subsubsection{\XeLaTeX}
% \label{sec:xelatex}
% 很多用户对 \LaTeX\ 命令执行的次数不太清楚。一个基本的原则是多次运行 \LaTeX\ 命
% 令直至不再出现警告。下面给出生成示例文档的详细过程（\texttt{\#} 开头的行为注
% 释），首先来看推荐的 \texttt{xelatex} 方式：
% \begin{shell}
% # 1. 发现里面的引用关系，文件后缀 .tex 可以省略
% $ xelatex thesis
%
% # 2. 编译参考文件源文件，生成 bbl 文件
% $ bibtex thesis
%
% # 3. 下面解决引用
% $ xelatex thesis
% $ xelatex thesis   # 如果不需要生成索引此时生成完整的 pdf 文件
% $ splitindex thesis -- -s hithesis.ist  # 自动生成索引
% $ xelatex thesis.tex
% \end{shell}
%
% \subsubsection{latexmk}
% \label{sec:latexmk}
% \texttt{latexmk} 命令支持全自动生成 \LaTeX\ 编写的文档，并且支持使用不同的工具
% 链来进行生成，它会自动运行多次工具直到交叉引用都被解决。下面给出了一个用
% \texttt{latexmk} 调用 \texttt{xelatex} 生成最终文档的示例：
% \begin{shell}
% # 一句话就够了！
% $ latexmk -xelatex thesis
% \end{shell}
%
% \subsubsection{make}
% \label{sec:make}
% \note[提示：]{若要使用 \texttt{make} 编译，需自行下载模板。因为 \TeX\ 发行版中
% 的 \file{Makefile} 不在当前目录。}
%
% 上面的方法虽然不复杂，但是每次都输入还是非常罗嗦，所以 \hithesis\ 提供了一
% 个 \file{Makefile}：
%
% 仓库根目录下的 \file{Makefile} 负责生成类文件与用户手册：
% \begin{shell}
% $ make clean
% $ make cls       # 生成 hithesisbook.cls 与 hithesisbook.cfg
% $ make doc       # 生成说明文档 hithesis.pdf
% \end{shell}
% 各示例文件夹（如 \file{examples/hitbook/chinese/}）下另有一份 \file{Makefile}，用于编译示例文档：
% \begin{shell}
% $ make thesis    # hitbook 示例生成 thesis.pdf；hitart 示例对应目标为 make report
% \end{shell}
%
% \hithesis\ 的 \file{Makefile} 默认用 \texttt{latexmk} 调用\texttt{xelatex} 编
% 译，此外还支持直接用 \texttt{xelatex} 编译。如有需要可修
% 改 \file{Makefile} 开头的参数或通过命令行传递参数（请参看 \file{README.md}），
% 进一步还可以修改 \file{latexmkrc} 进行定制。
%
% \subsection{升级}
% \label{sec:updgrade}
% \hithesis\ 升级非常简单，可以通过 \TeX 发行版的包管理工具自动更新发行版，也可
% 以下载最新的开发版，
% 将 \file{hithesis.ins}，\file{hithesis.dtx} 拷贝至工作目
% 录覆盖相应的文件，然后运行：
% \begin{shell}
% $ latex hithesis.ins
% \end{shell}
%
% 生成新的类文件和配置文件即可。也可以直接拷
% 贝 \file{hithesisbook.cls}、\file{hithesisbook.cfg}
% 与 \file{hithesis.ist}，免去上面命令的执行。
%
%
% \section{使用说明}
% \label{sec:usage}
% 本手册假定用户已经能处理一般的 \LaTeX\ 文档，并对 \BibTeX\ 有一定了解。如果
% 从来没有接触过 \TeX\ 和 \LaTeX，建议先学习相关的基础知识。
%
% \subsection{关于提问}
% \label{sec:howtoask}
% 按照优先级推荐提问的位置如下（三个 QQ 群不分先后）：
%
% \begin{itemize}
% \item \href{http://github.com/hithesis/hithesis/issues}{Github Issues}
% ^^A \item \href{https://git.oschina.net/dustincys/hithesis/issues}{OSChina Issues}
% \item hithesis讨论区：259959600 （人满）
% \item 窝工山hithesis派：851792460 （人满）
% \item HiThesis 分盟：704864357
% \end{itemize}
%
% \subsection{示例文件}
% \label{sec:userguide}
% 模板核心文件有三
% 项：\file{hithesisbook.cls}、\file{hithesisbook.cfg}
% 与 \file{hithesis.bst}，但是如果
% 没有示例文档用户会发现很难下手。所以推荐新用户从模板自带的示例文档入手，里面包
% 括了论文写作用到的所有命令及其使用方法，只需要用自己的内容进行相应替换就可以。
% 对于不清楚的命令可以查阅本手册。下面的例子描述了模板中章节的组织形式，来自于示
% 例文档，具体内容可以参考模板附带的 \file{thesis.tex}。
%
% \lstinputlisting[style=lstStyleLaTeX]{examples/hitbook/english/thesis.tex}
% \lstinputlisting[style=lstStyleLaTeX]{examples/hitbook/chinese/thesis.tex}
% \lstinputlisting[style=lstStyleLaTeX]{examples/hitart/reports/report.tex}
% \lstinputlisting[style=lstStyleLaTeX]{examples/hitart/reportplus/report.tex}
%
% \subsection{论文选项}
% \label{sec:option}
%
% 论文选项，就是在 \file{thesis.tex} 文件的开头，非注释的第一行的方括号中填写的选项，示例见上节。
% 各个选项的含义说明已经在上节中说明，所以这里就不重复了。
%
% \subsection{中文字体}
% \label{sec:chinese-fonts}
% 正确配置中文字体是使用模板的第一步。模板调用 \CTeX\ 宏包，只提供基于
% \pkg{xeCJK} 包，使用 \XeLaTeX\ 编译的方式。
% 关于如何使用字体命令、字号等等，属于模板格式范畴，在实现细节中讨论。
% 关于中文字体安装、配置的所有问题不在本模板讨论 范围。
%
% \subsection{前文}
% \label{sec:titlepage}
% 前文内容是正文之前，含封面、摘要、目录、符号表。
% 封面信息提供两种配置方法：一是通过统一设置命 令 \cs{hitsetup}
% 通过\emph{key=value} 形式完成；二是每个信息利用命令独立设置， 其中命令的名字跟
% \emph{key} 相同。两种方式可以交叉使用，并按顺序执行（即后来的设置会覆
% 盖前面的）。以 \texttt{c} 开头的命令跟中文相关，\texttt{e}
% 开头则为对应的英文。
%
% \DescribeMacro{\hitsetup}
% \cs{hitsetup} 用法与常见 \emph{key=value} 命令相同，如下：
% \begin{latex}
% \hitsetup{
%   key1 = value1,
%   key2 = {a value, with comma},
% }
% % 可以多次调用
% \hitsetup{
%   key3 = value3,
%   key1 = value11, % 覆盖 value1
% }
% \end{latex}
%
% \note[注意：]{\cs{hitsetup} 使用 \pkg{kvoptions} 机制，所以配置项之间不能有空行，否则
% 会报错。}
%
% 大多数命令的使用方法都是： \cs{command}\marg{arg}，例外者将具体指出。这些命令都
% 在示例文档的 \file{front/cover.tex} 中。
%
% \subsubsection{密级}
% \label{sec:setup-secret}
% \DescribeMacro{statesecrets}
% \DescribeMacro{natclassifiedindex}
% \DescribeMacro{intclassifiedindex}
% 定义秘密级别和国内国际索引号。
% \begin{latex}
% \hitsetup{
% statesecrets={公开},
% natclassifiedindex={TM301.2},
% intclassifiedindex={62-5},
% }
% \end{latex}
%
% \subsubsection{论文标题}
% \myentry{论文标题}
% \DescribeMacro{ctitle}
% \DescribeMacro{etitle}
% \DescribeMacro{ctitlecover}
% \DescribeMacro{ctitleone}
% \DescribeMacro{ctitletwo}
% \DescribeMacro{csubtitle}
% \DescribeMacro{esubtitle}
% 中英文标题。其中 \cs{ctitlecover} 仅用于封面，可在过长时手动断行；
% \cs{ctitleone}、\cs{ctitletwo} 仅用于本科生封面标题分行。
% 如果有副标题，需要在封面选项中设置subtitle=true，否则不显示副标题。
% \begin{latex}
% \hitsetup{
%   ctitlecover={封面中文题目可断行},
%   ctitle={论文中文题目},
%   etitle={Thesis English Title},
%   csubtitle={论文中文副题目（如果有）},
%   esubtitle={Thesis English Sub-Title (if necessary)},
%   ctitleone={本科生论文中文题目上部分},
%   ctitletwo={本科生论文中文题目下部分},
% }
% \end{latex}
%
% \subsubsection{作者姓名}
% \myentry{作者姓名}
% \DescribeMacro{cauthor}
% \DescribeMacro{eauthor}
% 作者姓名。
% \begin{latex}
% \hitsetup{
%   cauthor={中文姓名},
%   eauthor={Name in Pinyin}
% }
% \end{latex}
%
% \subsubsection{申请学位名称}
% \label{sec:degree}
% \myentry{学科名称}
% \DescribeMacro{cxueke}
% \DescribeMacro{exueke}
% 按照入学的培养计划中学科自行填写，具体学科名称不是本文档范畴。
% 学位名称（博士/硕士/学士）由文档类的 \option{type} 选项自动确定，无须手动设置。
%
% \begin{latex}
% \hitsetup{
%   cxueke={工学},
%   exueke={Engineering},
% }
% \end{latex}
%
% \subsubsection{院系名称}
% \myentry{院系名称}
% \DescribeMacro{caffil}
% \DescribeMacro{eaffil}
% 院系名称，同上，按照入学的培养计划中学科自行填写，具体院系名称不是本文档范畴。
% \begin{latex}
% \hitsetup{
%  caffil={机电工程学院},
%  eaffil={\emultiline[t]{School of Mechatronics Engineering \\ Mechatronics Engineering}},
% }
% \end{latex}
% \note[注意：]{个别学院英文名过长，使用以上方法自行换行。}
%
% \subsubsection{专业名称}
% \myentry{专业名称}
% \DescribeMacro{csubject}
% \DescribeMacro{esubject}
% 专业名称，同上，按照入学的培养计划中学科自行填写，具体名称不是本文档范畴。
% \begin{latex}
% \hitsetup{
%  csubject={机械制造及其自动化},
%  esubject={Computer Science and Technology},
% }
% \end{latex}
%
% \subsubsection{导师}
% \myentry{导师}
% \DescribeMacro{csupervisor}
% \DescribeMacro{esupervisor}
% 直接导师。
% \begin{latex}
% \hitsetup{
%   csupervisor={导师~教授},
%   esupervisor={Supervisor}
% }
% \end{latex}
%
% \myentry{副导师}
% \DescribeMacro{cassosupervisor}
% \DescribeMacro{eassosupervisor}
% 副指导教师。
% \begin{latex}
% \hitsetup{
%   cassosupervisor={副导师~副教授},
%   eassosupervisor={2nd Boss}
% }
% \end{latex}
%
% \myentry{联合导师}
% \DescribeMacro{ccosupervisor}
% \DescribeMacro{ecosupervisor}
% 硕士生联合指导教师，博士生联合导师。
% \begin{latex}
% \hitsetup{
%   ccosupervisor={联合导师~教授},
%   ecosupervisor={3rd Boss}
% }
% \end{latex}
%
% \subsubsection{成文日期}
% \myentry{日期}
% \DescribeMacro{cdate}
% \DescribeMacro{edate}
% 默认为当前时间，也可以自己指定。
% \begin{latex}
% \hitsetup{
%   cdate={中文日期},
%   edate={English Date},
% }
% \end{latex}
%
% \myentry{博士后日期}
% \DescribeMacro{cstartdate}
% \DescribeMacro{cenddate}
% \DescribeMacro{csubmitdate}
% \DescribeMacro{cfinishdate}
% 博士后出站报告封面需要填写进站、出站、提交、完成等多个日期，对应命令分别为
% \cs{cstartdate}、\cs{cenddate}、\cs{csubmitdate}、\cs{cfinishdate}。
% 其余学生类型的封面无须填写。
%
% \subsubsection{学生类型}
% \myentry{学生类型}
% \DescribeMacro{cstudenttype}
% \DescribeMacro{estudenttype}
% 非全日制教育申请学位者
%（同等学力人员）、（工程硕士）、（工商管理硕士）、
%（高级管理人员工商管理硕士）、（公共管理硕士）、（中职教师）、（高校教师）等，
% 具体要求按照入学的培养计划中学科自行填写，具体名称不是本文档范畴。
% \begin{latex}
% \hitsetup{
%   cstudenttype={同等学力人员},
%   estudenttype={Master of Art},
% }
% \end{latex}
%
% \subsubsection{学号}
% \myentry{学号}
% \DescribeMacro{cstudentid}
% 学号，具体要求按照入学的培养计划中学科自行填写。
% \begin{latex}
% \hitsetup{
%   cstudentid={9527},
% }
% \end{latex}
%
% \subsubsection{摘要}
% \myentry{摘要正文}
% \DescribeEnv{cabstract}
% \DescribeEnv{eabstract}
% \note[说明：]{摘要正文只能用环境命令的形式，不支持 \cs{hitsetup}。}
%
% \begin{latex}
% \begin{cabstract}
%  摘要请写在这里...
% \end{cabstract}
%
% \begin{eabstract}
%  Here comes the abstract in English...
% \end{eabstract}
% \end{latex}
%
% \myentry{关键词}
% \DescribeMacro{ckeywords}
% \DescribeMacro{ekeywords}
% 关键词用英文逗号分割写入相应的命令中，模板会解析各关键词并生成符合不同论文格式
% 要求的关键词格式。
% \begin{latex}
% \hitsetup{
%   ckeywords={关键词 1, 关键词 2},
%   ekeywords={keyword 1, keyword 2}
% }
% \end{latex}
%
% \subsubsection{符号对照表}
% \DescribeEnv{denotation}
% 主要符号表环境，单独在文件 \file{front/denotation.tex} 中。
% 跟据\PGR\
% 示例中要求，我工符号表是table环境，示例文件如下，由于我工要求博士论文图表标题是双语，所以任何对单个标题的全局格式调整都会影响到双语标题，所以这里使用\cs{vspace}，具体见实现细节中的描述。
% \begin{latex}
% \begin{denotation}
% \begin{table}[h]%此处最好是h
% \caption{国际单位制中具有专门名称的导出单位}
% \vspace{0.5em}\centering\wuhao
% \begin{tabular}{ccccc}
% \toprule
% 量的名称&单位名称&单位符号&其它表示实例\\
% \midrule
% 频率&赫[兹]&Hz&s-1\\
% \bottomrule
% \end{tabular}
% \end{table}
% \end{denotation}
% \end{latex}
%
% \subsubsection{目录}
% 目录不需要用户干预，自动生成，具体命令已经写在 \file{thesis.tex} 中。
%
% \subsection{正文}
%
% \subsubsection{图和表}
% \hit\ 博士毕业论文要求使用中英双语图题、表题，这增加了维护难度。
% 因为现有唯一的方法是在已有的图题或表题的基础上再添加一行英语图题或表题。
% 两个题之间的距离具体多少不在\PGR\ 中要求。目前的方法是用户手动调节该距离。
% 关于图题\PGR\ 和\UGR\ 只规定了居中，并没有规定居中对齐。然而评审老师很多喜欢居
% 中且居中对齐。模板默认选项是居中但不居中对齐，如果需要居中对齐，那么需要在
% \file{thesis.tex}的文档类选项中设置选项capcenterlast=true。详细方法见前文的介绍
% 。
% \begin{hitrgu}
% 每个图均应有图题（由图序和图名组成），图题不宜有标点符号，图名在图序之后空1个
% 半角字符排写。图序按章编排，如第1章第一个插图的图号为“图1-1”。图题置于图下，硕
% 士论文只用中文，博士论文用中、英两种文字，居中书写，中文在上，要求中文用宋体5
% 号字，英文用Times New Roman 5号字。有图注或其它说明时应置于图题之上。引用图应
% 注明出处，在图题右上角加引用文献号。图中若有分图时，分图题置于分图之下或图题之
% 下，可以只用中文书写，分图号用a)、b)等表示。图中各部分说明应采用中文（引用的外
% 文图除外）或数字符号，各项文字说明置于图题之上（有分图时，置于分图题之上）。图
% 中文字用宋体、Times New Roman字体，字号尽量采用5号字（当字数较多时可用小5号字
% ，以清晰表达为原则，但在一个插图内字号要统一）。同一图内使用文字应统一。图表中
% 物理量、符号用斜体。
% \end{hitrgu}
% 单双语图题的方法如下，注释中说明。
% \begin{latex}
% \begin{figure}[htpb]
% \centering
% \includegraphics[width = 0.4\textwidth]{golfer}
% \bicaption[golfer1]{}{注意图中文字尽量用五号字
% }{Fig.$\!$}{The person playing golf}
% \end{figure}
% \end{latex}
% 单张单图题的格式如下，
% \begin{latex}
% \begin{figure}[h]
% \centering
% \includegraphics[width = 0.4\textwidth]{golfer}
% \caption{注意图中文字字号尽量用五号字}
% \end{figure}
% \end{latex}
% 并排图例。
% \begin{latex}
% \begin{figure}[htbp]
% \centering
% \begin{minipage}{0.4\textwidth}
% \centering
% \includegraphics[width=\textwidth]{golfer}
% \bicaption[golfer2]{}{打高尔夫球的人}{Fig.$\!$}{The person playing golf}
% \end{minipage}
% \begin{minipage}{0.4\textwidth}
% \centering
% \includegraphics[width=\textwidth]{golfer}
% \bicaption[golfer3]{}{打高尔夫球的人}{Fig.$\!$}{The person playing golf}
% \end{minipage}
% \end{figure}
% \end{latex}
% 子图图例。
% \begin{latex}
% \begin{figure}[htbp]
% \centering
% \subfigure{\label{golfer41}}\addtocounter{subfigure}{-2}
% \subfigure[The person playing golf]{\subfigure[打高尔夫球的人~1]{\includegraphics[width=0.4\textwidth]{golfer}}}
% \subfigure{\label{golfer42}}\addtocounter{subfigure}{-2}
% \subfigure[The person playing golf]{\subfigure[打高尔夫球的人~2]{\includegraphics[width=0.4\textwidth]{golfer}}}
% \bicaption[golfer4]{}{打高尔夫球的人}{Fig.$\!$}{The person playing golf}
% \end{figure}
% \end{latex}
% 表格示例，表格中的字体是可以自行调整的。
% \begin{latex}
% \begin{table}[htbp]
% \bicaption[table1]{}{符合研究生院绘图规范的表格}{Table$\!$}{Table in agreement of the standard from graduate school}
% \vspace{0.5em}\centering\wuhao
% \begin{tabular}{ccccc}
% \toprule
% $D$(in) & $P_u$(lbs) & $u_u$(in) & $\beta$ & $G_f$(psi.in)\\
% \midrule
%  5 & 269.8 & 0.000674 & 1.79 & 0.04089\\
% 10 & 421.0 & 0.001035 & 3.59 & 0.04089\\
% 20 & 640.2 & 0.001565 & 7.18 & 0.04089\\
% \bottomrule
% \end{tabular}
% \end{table}
% \end{latex}
% 因为长表格不是浮动体，不会自动调整位置、也不会自动调整字体大小，一切都要手动设
% 置。特别繁琐。
% \begin{latex}
% \ltfontsize{\dawu[1.667]} %设置表格内字体行间距
% \dawu[1.667]\begin{longtable}{ccc} % 注意此处设置的是表格线距离
% \longbionenumcaption{}{{\wuhao 中国省级行政单位一览 %此处要添加字体设置
% }\label{table2}}{Table$\!$}{}{{\wuhao Overview of the provincial administrative
% unit of China}}{-0.5em}{3.15bp}\\ %注意后两个参数分别是中英标题间距、标题和表格的间距。
% %\caption{\wuhao 中国省级行政单位一览}\\[1em] %注意此处是标题和表格间距，这行
% %是单语标题
% \toprule 名称 & 简称 & 省会或首府  \\ \midrule
% \endfirsthead
% \multicolumn{3}{r}{表~\thetable（续表）}\vspace{0.5em}\\
% \toprule 名称 & 简称 & 省会或首府  \\ \midrule
% \endhead
% \bottomrule
% \endfoot
% 北京市 & 京 & 北京\\
% 天津市 & 津 & 天津\\
% 河北省 & 冀 & 石家庄市\\
% 山西省 & 晋 & 太原市\\
% 内蒙古自治区 & 蒙 & 呼和浩特市\\
% 辽宁省 & 辽 & 沈阳市\\
% 吉林省 & 吉 & 长春市\\
% 黑龙江省 & 黑 & 哈尔滨市\\
% 上海市 & 沪/申 & 上海\\
% 江苏省 & 苏 & 南京市\\
% 浙江省 & 浙 & 杭州市\\
% 安徽省 & 皖 & 合肥市\\
% 福建省 & 闽 & 福州市\\
% 江西省 & 赣 & 南昌市\\
% 山东省 & 鲁 & 济南市\\
% 河南省 & 豫 & 郑州市\\
% 湖北省 & 鄂 & 武汉市\\
% 湖南省 & 湘 & 长沙市\\
% 广东省 & 粤 & 广州市\\
% 广西壮族自治区 & 桂 & 南宁市\\
% 海南省 & 琼 & 海口市\\
% 重庆市 & 渝 & 重庆\\
% 四川省 & 川/蜀 & 成都市\\
% 贵州省 & 黔/贵 & 贵阳市\\
% 云南省 & 云/滇 & 昆明市\\
% 西藏自治区 & 藏 & 拉萨市\\
% 陕西省 & 陕/秦 & 西安市\\
% 甘肃省 & 甘/陇 & 兰州市\\
% 青海省 & 青 & 西宁市\\
% 宁夏回族自治区 & 宁 & 银川市\\
% 新疆维吾尔自治区 & 新 & 乌鲁木齐市\\
% 香港特别行政区 & 港 & 香港\\
% 澳门特别行政区 & 澳 & 澳门\\
% 台湾省 & 台 & 台北市\\
% \end{longtable}\normalsize %注意这里要恢复正常字体
% \end{latex}
% \subsubsection{公式}
% 公式不做介绍，与正常用法一致。
% \subsubsection{数学环境}
% \label{sec:math}
% \hithesis\ 定义了常用的数学环境：
%
% \begin{center}
% \begin{tabular}{*{7}{l}}\toprule
%   axiom & theorem & definition & proposition & lemma & conjecture &\\
%   公理 & 定理 & 定义 & 命题 & 引理 & 猜想 &\\\midrule
%   proof & corollary & example & exercise & assumption & remark & problem \\
%   证明 & 推论 & 例子& 练习 & 假设 & 注释 & 问题\\\bottomrule
% \end{tabular}
% \end{center}
%
% 比如：
% \begin{latex}
% \begin{definition}
%   道千乘之国，敬事而信，节用而爱人，使民以时。
% \end{definition}
% \end{latex}
% 产生（自动编号）：
% \medskip
%
% \noindent\framebox[\linewidth][l]{{\heiti 定义~1.1~~~} % {道千乘之国，敬事而信，节用而爱人，使民以时。}}
%
% \smallskip
% 列举出来的数学环境毕竟是有限的，如果想用\emph{胡说}这样的数学环境，那么可以定义：
% \begin{latex}
% \newtheorem{nonsense}{胡说}[chapter]
% \end{latex}
%
% 然后这样使用：
% \begin{latex}
% \begin{nonsense}
%   契丹武士要来中原夺武林秘笈。—— 慕容博
% \end{nonsense}
% \end{latex}
% 产生（自动编号）：
%
% \medskip
% \noindent\framebox[\linewidth][l]{{\heiti 胡说~1.1~~~} % {契丹武士要来中原夺武林秘笈。—— 慕容博}}
% \subsubsection{算法}
% 我工算法不在规范中要求且一千个评审老师有一千个算法格式喜好。详见
% \href{https://github.com/PlutoThesis/PlutoThesis#%E6%B2%A1%E6%9C%89%E6%98%8E%E7%A1%AE%E8%A6%81%E6%B1%82%E7%9A%84%E6%A0%BC%E5%BC%8F}{PlutoThesis}
% 中的各个实验室算法喜好举例。在此多说无益。
% \subsubsection{引用参考文献}
% \DescribeMacro{\inlinecite}
% 学校要求的参考文献引用有两种模式：（1）上标模式。比如``同样的工作有很
% 多$^{[1,2]}$\ldots''。（2）正文模式。比如``文[3] 中详细说明了\ldots''。其中上标
% 模式使用远比正文模式频繁，所以为了符合使用习惯，上标模式仍然用常规
% 的 \cs{cite}\marg{key}，而 \cs{inlinecite}\marg{key} 则用来生成正文模式。
%
% 关于参考文献模板推荐使用 \BibTeX，关于中文参考文献需要额外增加一个 Entry:
% \texttt{language}，将其设置为 \texttt{zh} 用来指示此参考文献为中文，以
% 便 \file{hithesis.bst} 处理。如：
% \begin{latex}
% @INPROCEEDINGS{cnproceed,
%   author    = {王重阳 and 黄药师 and 欧阳峰 and 洪七公 and 段皇帝},
%   title     = {武林高手从入门到精通},
%   booktitle = {第~$N$~次华山论剑},
%   year      = 2006,
%   address   = {西安, 中国},
%   language      = "zh",
% }
%
% @ARTICLE{cnarticle,
%   AUTHOR  = "贾宝玉 and 林黛玉 and 薛宝钗 and 贾探春",
%   TITLE   = "论刘姥姥食量大如牛之现实意义",
%   JOURNAL = "红楼梦杂谈",
%   PAGES   = "260--266",
%   VOLUME  = "224",
%   YEAR    = "1800",
%   LANGUAGE    = "zh",
% }
% \end{latex}
%
% 注意如果不需要引用参考文献，请删除 \file{thesis.tex} 中 \cs{bibliography} 开头的两行，
% 以避免可能的编译错误。
%
% \subsubsection{列表环境}
% \DescribeEnv{itemize}
% \DescribeEnv{enumerate}
% \DescribeEnv{description}
% 为了适合中文习惯，模板将这三个常用的列表环境用 \pkg{enumitem} 进行了纵向间距压
% 缩。一方面清除了多余空间，另一方面用户可以自己指定列表环境的样式（如标签符号，
% 缩进等）。细节请参看 \pkg{enumitem} 文档，此处不再赘述。
% \subsection{后文}
%
% \subsubsection{结论}
% \DescribeEnv{conclusions}
% 结论之后为后文内容。
% \begin{latex}
% \begin{conclusions}
%
% 学位论文的结论作为论文正文的最后一章单独排写，但不加章标题序号。
%
% 结论应是作者在学位论文研究过程中所取得的创新性成果的概要总结，不能与摘要混为一
% 谈。博士学位论文结论应包括论文的主要结果、创新点、展望三部分，在结论中应概括论
% 文的核心观点，明确、客观地指出本研究内容的创新性成果（含新见解、新观点、方法创
% 新、技术创新、理论创新），并指出今后进一步在本研究方向进行研究工作的展望与设想
% 。对所取得的创新性成果应注意从定性和定量两方面给出科学、准确的评价，分（1）、
% （2）、（3）…条列出，宜用“提出了”、“建立了”等词叙述。
%
% \end{conclusions}
% \end{latex}
%
% \subsubsection{参考文献}
% 在后文中的参考文献是自动生成的，不需要用户干预，具体命令在 \file{thesis.tex} 中有
% 示例。
%
% \subsubsection{附录}
% \DescribeEnv{appendix}
% 所有的附录都插到这里来。因为附录会更改默认的 chapter 属性，而后面的{\heiti 个人简
%   历}又需要恢复，所以实现为环境可以保证全局的属性不受影响。
% \begin{latex}
% \begin{appendix}
% \input{back/appA.tex}
% \end{appendix}
% \end{latex}
%
% \subsubsection{所发表文章}
% \DescribeEnv{publication}
% 虽然在\PGR\UGR\ 中都没有明确规定此处的格式，但按照旧模板PlutoThesis，此处格式
% 非常复杂。此处仍然使用旧模板中的设置方法。
% \begin{latex}
% \begin{publication}
% \noindent\textbf{（一）发表的学术论文}
% \begin{publist}
% \item	XXX，XXX. Static Oxidation Model of Al-Mg/C Dissipation Thermal Protection Materials[J]. Rare Metal Materials and Engineering, 2010, 39(Suppl. 1): 520-524.（SCI~收录，IDS号为~669JS，IF=0.16）
% \item XXX，XXX. 精密超声振动切削单晶铜的计算机仿真研究[J]. 系统仿真学报，2007，19（4）：738-741，753.（EI~收录号：20071310514841）
% \item XXX，XXX. 局部多孔质气体静压轴向轴承静态特性的数值求解[J]. 摩擦学学报，2007（1）：68-72.（EI~收录号：20071510544816）
% \item XXX，XXX. 硬脆光学晶体材料超精密切削理论研究综述[J]. 机械工程学报，2003，39（8）：15-22.（EI~收录号：2004088028875）
% \item XXX，XXX. 基于遗传算法的超精密切削加工表面粗糙度预测模型的参数辨识以及切削参数优化[J]. 机械工程学报，2005，41（11）：158-162.（EI~收录号：2006039650087）
% \item XXX，XXX. Discrete Sliding Mode Cintrok with Fuzzy Adaptive Reaching Law on 6-PEES Parallel Robot[C]. Intelligent System Design and Applications, Jinan, 2006: 649-652.（EI~收录号：20073210746529）
% \end{publist}
%
% \noindent\textbf{（二）申请及已获得的专利（无专利时此项不必列出）}
% \begin{publist}
% \item XXX，XXX. 一种温热外敷药制备方案：中国，88105607.3[P]. 1989-07-26.
% \end{publist}
%
% \noindent\textbf{（三）参与的科研项目及获奖情况}
% \begin{publist}
% \item	XXX，XXX. XX~气体静压轴承技术研究, XX~省自然科学基金项目.课题编号：XXXX.
% \item XXX，XXX. XX~静载下预应力混凝土房屋结构设计统一理论. 黑江省科学技术二等奖, 2007.
% \end{publist}
% %\vfill
% %\hangafter=1\hangindent=2em\noindent
% %\setlength{\parindent}{2em}
% \end{publication}
% \end{latex}
%
% \subsubsection{索引}
% \DescribeEnv{ceindex}
% 我工要求中英文双语索引。后文中的自动索引实际上不需要用户干预。
%\begin{latex}
% \begin{ceindex}
%   %如果想要手动加索引，注释掉以下这一样，用wordlist环境
% \printsubindex*
% \end{ceindex}
%\end{latex}
% 手工添加索引的方法不推荐，模板中将去除该功能。
% \subsubsection{授权}
% \DescribeMacro{\authorization}
% 授权页中的签名和日期是需要手写，不需要人工干预。具体示例在 \file{thesis.tex} 中。
%\begin{latex}
% \authorization %授权
% %\authorization[scan.pdf] %添加扫描页的命令，与上互斥
%\end{latex}
%
% \subsubsection{致谢声明}
% \DescribeEnv{acknowledgement}
% 把致谢做成一个环境更好一些，直接往里面写感谢的话就可以啦！
%
% \begin{latex}
% \begin{acknowledgement}
%   …
%   感谢\hit\LaTeX\ 论文模板\hithesis\ !
% \end{acknowledgement}
% \end{latex}
%
%
% \subsubsection{简历}
% \DescribeEnv{resume}
% 个人简历。
% 实际上，致谢和个人简历是自由发挥的地区，字体，文体，格式，内容，完全自己决定。
% \begin{latex}
% \begin{resume}
% XXXX~年~XX~月~XX~日出生于~XXXX。
%
% XXXX~年~XX~月考入~XX~大学~XX~院（系）XX~专业，XXXX~年~XX~月本科毕业并获得~XX~学学士学位。
%
% XXXX~年~XX~月------XXXX~年~XX~月在~XX~大学~XX~院（系）XX~学科学习并获得~XX~学硕士学位。
%
% XXXX~年~XX~月------XXXX~年~XX~月在~XX~大学~XX~院（系）XX~学科学习并获得~XX~学博士学位。
%
% 获奖情况：如获三好学生、优秀团干部、X~奖学金等（不含科研学术获奖）。
%
% 工作经历：
% \end{resume}
% \end{latex}
%
% \subsection{其它}
% 模板的配置文件 \file{hithesisbook.cfg} 中定义了很多固定词汇，一般无须修改。如果有特殊需求，
% 推荐在导言区使用 \cs{renewcommand}。
%
%
% \subsection{捐助}
% \changes{v1.0.01}{2017/08/27}{添加了捐助、矢量化本科论文模板的图片logo}
% 各位刀客和大侠如用的嗨，要解囊相助，请支付宝赞助
% 参照图~\ref{zfb}~中提示操作（二维码被矢量化后之后去
% 除了头像等冗余无用的部分～）。
% \begin{figure}[h]
% \centering
\includegraphics[width=0.5\textwidth]{zfb.eps}
% \caption{支付宝}
% \label{zfb}
% \end{figure}
%
